% ===== §7 =====
\section{The Chinese \emph{Wen--Hua} Tradition}
\label{p24:sec:chinese}

\noindent
The five traditions examined so far are Western, and they share, for all their differences, a common descent from the social sciences and the philosophy of the modern West. There is a body of Chinese reflection on \emph{wen} (\zh{文}) and \emph{hua} (\zh{化}) that does not belong to that descent and is not, in the strict sense, a theory of culture in the modern disciplinary meaning at all: it is closer to a discourse on civilization, cultivation, and the forming of persons and a world. It is included here as an independent strand, and not merely as an epigraph, because its difference of concern preserves an insight the Western traditions, with the partial exception of the dialogic, tend to lose---and because the terms of the intimacy series have from the beginning drawn on this material. The section confines itself to what can be stated with confidence; where it would require a philological or exegetical judgement the author is better placed to make, it marks the place and leaves it open rather than filling it with a claim it cannot warrant.

\subsection*{What can be stated: the verbal force of \emph{hua}}

One observation can be made with confidence, and it is the one that matters most for the argument, because it is grammatical before it is interpretive. The modern compound \emph{wenhua} is built of two characters whose senses pull in the direction this paper has been arguing. \emph{Wen} is pattern, the woven grain, the mark or sign---the nominal, deposited face of culture, close to what the transmission tradition sees. \emph{Hua} is transformation, and its force is \emph{verbal}: to \emph{hua} is to transform, to convert, to bring to form, to cause to change or grow. Where the Latin \emph{cultura} leans, through its agricultural root, toward the tended growth that becomes a possession, \emph{hua} names the transforming itself, the ongoing action rather than its settled product. This is not an exotic subtlety but a plain feature of the term, and it is exactly the feature the settled Western traditions let fall: that at the heart of the very word for culture, in one of its two components, stands a verb of generation. The Chinese compound thus holds, in its internal tension between \emph{wen} and \emph{hua}, the same tension this paper has set at the centre of its ontology, between the pattern that is deposited and the transformation that generates it; and it holds the generative pole in a term that will not let it be forgotten.

The locus classicus for reading culture as an act of formation rather than a deposit is the phrase from the \emph{Tuan} commentary on the hexagram Bi (\zh{贲}) in the \emph{Book of Changes}---which may be rendered ``observe the patterns of the human, and by them bring the world to form'' \parencite{legge1963yijing}. What can be said with confidence, and only this, is that here \emph{huacheng} (the bringing-to-form) names an ongoing cultivation of a shared world, rather than the handing-on of a finished heritage: the character \emph{hua} does the work of a verb of formation, consistent with the grammatical point just made. Any more particular claim about the passage---its place within the philosophy of the \emph{Changes}, the commentarial tradition that reads it, its relation to the classical theory of \emph{liyue} (rites and music), and the exact weight of \emph{renwen} (the human pattern) in it---is a matter of exegesis the present author is better placed to supply, and it is left open here rather than asserted.

\subsection*{Marked and left open}

\noindent\emph{The following are deliberately marked and left open, to be supplied or verified by the author, who has the classical competence the argument here requires; the section does not assert them, and states only the grammatical point above, which it can warrant.}
\begin{itemize}
  \item The Confucian account of \emph{liyue} as a cultivation that forms persons and generates a relational order, and its bearing on the thesis that culture is generative rather than deposited --- to be developed with the specific textual support the author judges apt.
  \item The Daoist material on \emph{daosheng} (the Ways begetting) and \emph{ziran} (self-so, spontaneity)---generation as spontaneous, non-coercive bringing-forth---and its connection to the treatment of water and of \emph{xuande} (dark virtue) in the earlier papers of this series, stated only to the extent the author wishes to warrant.
  \item The exact sense in which \emph{hua} is to be distinguished from mere change, and whether a stronger claim than the grammatical one above can be responsibly made.
\end{itemize}

\subsection*{The strand's function, and its internal tension}

Whatever the detail with which the classical material is finally filled in, the strand's function in the argument can be stated now, for it is structural. It supplies, from outside the Western descent, the one thing that descent tends to lose: a conception of culture in which the generative, verbal, ongoing character of formation is not an afterthought but is lodged in the very term. In the map of the second section it therefore falls at the generative pole, near the dialogic, and it is a second witness---independent of Bakhtin, and older---to the truth that culture is process before it is product.

But the strand carries a tension that must be held, and it is the same dialectical caution the paper has applied to every tradition. The forming of a world and the cultivation of persons have, historically, also been instruments of a hierarchical order: a \emph{huacheng} directed from above, a cultivation that forms subjects to their places in a structure of authority. The generative insight and the pedagogy of hierarchy have coexisted in this tradition, and to take the insight without noticing the hierarchy would be to commit, on Chinese material, exactly the naivety the power tradition of the previous section warned against. The strand is therefore not a refuge from the problem of power but another instance of it, to be evaluated dialectically: its conception of generative formation is a real gift to the argument, and its historical entanglement with a cultivation imposed from above is a real limit, and both are carried forward---the gift to the synthesis, the limit to the treatment of power and of cultural evolution that the synthesis will require.
